\section{Character Creation}\label{Character Creation}

  Creating a character involves a mixture of thematic and mechanical decisions that will work together to create a fun character that is rewarding to play.
  As mentioned earlier in this chapter, there are four core systems for customizing your character's mechanics: class, attributes, skills, and species.
  In addition, there are five core thematic considerations when creating a character: concept, personality, motivation, background, and appearance.

  These decisions are described below in an order that makes sense for many characters, and full details for each decision are given after this initial list.
  It is essentially a sandwich, with narrative decisions wrapped around a central core of your character's mechanical components.
  However, you can make several of these decisions in any order, and you may find it easier to create a character in a different way.
  The only real limitation is that your skills should generally be the last mechanical choice you make, since they are strongly affected by your class and attributes.

  \begin{enumerate}
    \item Character concept: Describe your character with a short, simple phrase that captures their essence.
    \item Motivation and goal: Describe what your character wants.

    \item Species: Define your character's species.
    \item Attributes: Define your character's fundamental physical and mental potential.
    \item Class archetypes: Define your character's source of power.
    \item Insight points: Learn new abilities.
    \item Skills: Define your character's areas of non-combat expertise.

    \item Personality: Describe how your character acts and reacts to the world.
    \item Background: Describe what made your character become who they are now.
    \item Appearance: Describe what your character looks like.
    \item Alignment: Describe your character's moral compass.
    \item Name: Choose a name.
  \end{enumerate}

  \subsection{Step 1: Character Concept}

    Fundamentally, who is your character?
    You should think of a short phrase that describes the core concept behind the character you will create.
    It's best to think in broad strokes when creating a character concept.
    Your concept should be more than just a factual description of your species or what you do.
    It should be something that makes you memorable.
    Some sample character concepts are given below for inspiration.

    \begin{raggeditemize}
      \item Pragmatic wanderer
      \item Artistic pixie
      \item Mushroom-obsessed hermit
      \item Bumbling do-gooder
      \item Dim-witted bodyguard
      \item Cowardly storyteller
      \item Bear-barian
      \item Parsimonious law enforcer
      \item Peaceful naturalist
      \item Trigger-happy pyromaniac
      \item Heroic, simple-minded warrior
      \item Friendly necromancer
      \item Chaotic speed demon
      \item Pompous ex-noble
      \item Sarcastic mercenary
      \item Battle-scarred priest
      \item Ambitious arcane prodigy
      \item Charismatic musician
      \item Aloof scholar
      \item Blunt-spoken warrior
      \item Crazed prophet
      \item Polite warrior
      \item World-weary pirate
      \item Devout cultist
      \item Con artist with a heart of gold
    \end{raggeditemize}

  \subsection{Step 2: Motivation and Goal}
    Why does your character put in all of the effort that adventuring requires?
    They probably have a goal that they are trying to achieve, or an ideal that they are trying to embody.
    Writing down a specific goal or ideal can be helpful as an anchor point when defining the character.

  \subsection{Step 3: Species}
    It's often convenient to make your species your first mechanically relevant choice.
    Your species can have a strong effect on your personality and narrative, but it has a relatively small effect on your character's play style.
    It's also easier to know your species before you choose your attributes, since your species can slightly modify your attributes.

    Choose one of the eight common species options, or talk with your GM about choosing an uncommon species (see \pcref{Uncommon Species}).
    Record any specific abilities the species gives you on your character sheet, but if this is your first mechanical choice, you won't be able to finalize any of your statistics yet.
    You should also choose the languages that you can speak, since that is influenced by your species (see \pcref{Communication and Languages}).

  \subsection{Step 4: Attributes}\label{Step 4: Attributes}
    Your attributes are a good option for your second mechanically relevant choice.
    They have a large impact on your character's strengths and weaknesses, so it's useful to know them as soon as possible.
    They're also much easier to understand and finalize than your class archetypes.

    You have 8 points to distribute among your attributes.
    Increasing an attribute by 1 costs 1 point, and you can increase each attribute up to a maximum of 3.
    Instead of allocating points yourself, you can use one of the following three common attribute arrays:
    \begin{raggeditemize}
      % 3 + 2 + 2 + 1 = 8
      \item Standard: 3, 2, 2, 1, 0, 0
        % 3 + 3 + 2 = 8
      \item Specialized: 3, 3, 2, 0, 0, 0
        % 2 + 2 + 2 + 1 + 1 = 8
      \item Balanced: 2, 2, 2, 1, 1, 0
    \end{raggeditemize}

    Once you have chosen your attributes, add your species modifier to your attributes (if any).
    Then, record in your character sheet the various effects that your attributes have on your statistics.

    \subsubsection{Attribute Penalties}\label{Attribute Penalties}
      You can voluntarily take penalties to your attributes.
      If you reduce an attribute by a total of \minus1, you gain an additional \glossterm{trained skill} (see \pcref{Trained Skills}).
      If you reduce an attribute by a total of \minus2, you instead gain an additional \glossterm{insight point} (see \pcref{Insight Points}).
      You cannot gain these benefits from reducing more than two attributes below 0 in this way.
      In addition, you can never reduce an attribute below \minus2 in this way.

  \subsection{Step 5: Class and Class Archetypes}
    This is the most complicated choice you have to make for your character.
    It requires reading at least some of the Classes chapter to understand which classes are interesting to you.
    Class details can be found in \pcref{Classes}.

    You should choose one of the eleven classes, and apply all effects of choosing that as your \glossterm{base class}.
    Then, choose one of the five archetypes within that class.
    You gain the rank 1 ability from that archetype (see \pcref{Archetypes}).

    When you reach levels 2 and 3, you'll choose new archetypes from the same class, becoming rank 1 in each of those archetypes as well.
    After that, you won't gain any more new archetypes when you gain levels.
    Instead, you'll just increase your rank in the three archetypes you already have.

    If you are particularly adventurous, this is also when you should choose if you want to be a multiclass character.
    Multiclass characters can gain archetypes from multiple classes.
    This does not increase the number of archetypes you know, so it does not directly increase your power.
    However, multiclass characters can be more specialized or more versatile than single-class characters, and can represent unusual character concepts.
    For details, see \pcref{Multiclass Characters}.

  \subsection{Step 6: Insight Points}
    Once you have chosen your class archetypes, attributes, and species, you know how many insight points you have, and can choose how to spend them.
    Don't forget to record on your character sheet how you spent each insight point.
    Otherwise, you might get confused later about why you have more spells known than you normally would.

    In some circumstances, you might want to delay spending your insight points until you are higher level.
    For example, a fighter/sorcerer multiclass character who wants to have both spells and maneuvers can't have access to both spells and maneuvers at level 1, so they wouldn't be able to spend insight points on both spells and maneuvers.
    You aren't forced to spend all of your insight points, so you can save them up for later.
    You can also talk to your GM about spending them at level 1 and then retraining those insight points once you are higher level.

  \subsection{Step 7: Skills}
    You should choose which skills you have as \glossterm{trained skills} (see \pcref{Skills}).
    Your \glossterm{class} gives you a certain number of trained skills from among its \glossterm{class skills}.
    The class skills for each class are summarized in \tref{Class Skills}.

    There are other ways to become trained in skills that are not part of your class.
    If your Intelligence is positive, you gain additional trained skills equal to your Intelligence.
    You can also spend \glossterm{insight points} to gain one trained skill per insight point (see \pcref{Insight Points}).
    Some abilities can grant additional trained skills.

    If you are untrained in a skill, your bonus with that skill is equal to half of its associated attribute (if any).
    If you are trained in a skill, your bonus with that skill is equal to 3 \add the higher of its associated attribute (if any) and half your level.
    Many abilities can increase or decrease your bonus with particular skills.

    The number of skills you can have trained, and which skills those are, depend on every preceding step, so it's a good place to finish.

    Sometimes, you might have more trained skills than you know what to do with, especially if you are still figuring out the details of your character concept.
    You aren't forced to decide all of your trained skills at level 1, so you can save them up and choose more trained skills when you level up.
    You can also talk to your GM about letting you decide your trained skills on the fly during the first game session or two based on what actions you take during the session.
    This can be a fun way to figure out what your character's personality is through the process of playing them.

  \subsection{Step 8: Starting Equipment}
    When you create a character, they can start with some basic items.
    Items have \glossterm{item ranks} that indicate the approximate rank that characters can reasonably get access to them.
    Typically, you can start with a single rank 1 item, up to three rank 0 items, and a standard adventuring kit.
    Individual campaigns or character backstories may significantly change what starting equipment is available, so check with your GM.

  \subsection{Step 9: Personality}

    How does your character behave?
    You should decide, in broad terms, what your character's personality is.
    This will change over time, especially as you start playing the character in the game, so you don't need to define everything perfectly.
    However, having a general sense of how your character behaves is helpful.

    For most games, it's important to have a personality that can tolerate working with others in a group.
    A character that is excessively aloof, moody, or obnoxious can make the game more difficult to enjoy for everyone.
    Likewise, a character who tries to speak for everyone or who repeatedly steals the spotlight from others can be frustrating to work with.
    You should figure out the right balance with your fellow players and your GM.\@

  \subsection{Step 10: Background}
    What happened in their character's past to make them the way that they are?
    What were their parents like, and where are they now?
    You don't have to have all of the answers when you first create a character, but it's good to have some idea.
    The richer your backstory, the more the GM can weave that into the narrative of the current story.
    Sometimes, it's fun to take a break from saving the world to go visit someone's grandma.
    For details about suggested backgrounds that have a strong effect on the game world, see \pcref{Backgrounds}.

  \subsection{Step 11: Appearance}
    What does your character look like?
    What would someone's first impression of them be?
    This can be helpful for understanding how other characters in the game world - or even monsters - would react to you.

  \subsection{Step 12: Alignment}
    Your character's alignment reflects their moral character: are they more inclined to good or evil, and to chaos or order?
    Alignments are described in more detail at \pcref{Alignment}.

  \subsection{Step 13: Name}
    What is your character's name?
    This choice can influence the tone your character will set in the game.
    If your name is Sir Patty Cakes or Shanky, the game is likely to be lighter and sillier in tone.
    Fancy fantasy-appropriate names like Ayala or Theodolus tend to push the game in a slightly more serious direction, especially if you make the daring choice to include a canonical last name.
    As always, stay in tune with what the GM and the other players are expecting.
